% Copyright (c) 2026 Geovana Neves. Licensed under CC BY 4.0 (https://creativecommons.org/licenses/by/4.0/). See LICENSE-AND-AUTHORSHIP.md.
%
% 00-fts-overview/text/02-theory.tex
%
% Part 2: the theory a user needs to justify the boundary-condition and setup
% choices, steady, unsteady and rotating, and nothing more. Manual pages are
% those of build 26121 [6]; the physics is in [7-12, 15-17].

\section{How a panel solver works}

\begin{frame}{Potential flow and its boundary conditions}
\begin{columns}[T]
\begin{column}{0.56\textwidth}
\begin{itemize}\small
  \item Assumptions: \textbf{incompressible, irrotational, inviscid}. The
        velocity derives from a potential, \(\mathbf{V} = \nabla\phi\), and
        mass conservation becomes the Laplace equation,
        \(\nabla^2\phi = 0\) [9, 10].
  \item The equation is \textbf{linear}: elementary solutions superpose, and
        the flow about a body is assembled from them.
  \item Two conditions close it: \textbf{flow tangency} on the surface, and
        \textbf{decay} of the disturbance far away.
\end{itemize}
\end{column}
\begin{column}{0.40\textwidth}
\begin{block}{What is absent}
\small No volume grid, so \textbf{no farfield to size} and no blockage to
correct: the disturbance decays because the elementary solutions decay
[6, p.~164].
\end{block}
\begin{block}{What it can predict}
\small Pressure effects: lift, induced drag, moments, interference. Never
friction drag, which lives in a boundary layer it does not model.
\end{block}
\end{column}
\end{columns}
\end{frame}

\begin{frame}{From source and doublet panels to surface vorticity}
\begin{columns}[T]
\begin{column}{0.58\textwidth}
\begin{itemize}\small
  \item A classic panel method gives each surface panel a source or doublet
        strength, enforces tangency at one point per panel, and solves a
        dense linear system [9, 10].
  \item FlightStream carries \textbf{vorticity} on the surface mesh instead
        [7, 8]. Degenerate faces are removed in preprocessing, and
        flow-aligned quadrilaterals are preferred on lifting surfaces
        [6, pp.~75, 156].
  \item Whatever the singularity: \textbf{surface resolution drives the
        induced velocities}, and cost grows steeply with the panel count.
\end{itemize}
\end{column}
\begin{column}{0.38\textwidth}
\slidefig[0.46]{g00-panels-and-wake}
\end{column}
\end{columns}
\begin{takeaway}Different singularity, same physics: enforce tangency, resolve
the circulation, integrate the pressures.\end{takeaway}
\end{frame}

\begin{frame}{The Kutta condition and the wake}
\begin{itemize}\small
  \item Potential flow alone \textbf{does not fix lift}: any circulation
        satisfies the equations. The physical statement that selects the real
        solution is smooth flow leaving a sharp trailing edge [10, 11].
  \item In the solver, \textbf{marking a trailing edge} is that statement:
        wake filaments are shed from marked edges, and an unmarked edge sheds
        nothing at all [6, p.~165].
  \item The wake carries the trailed vorticity downstream. Its strength
        mirrors the spanwise load gradient, and \textbf{induced drag} is the
        price of that vorticity.
\end{itemize}
\begin{alertblock}{The most critical input}
\small An unmarked trailing edge raises no error. It returns a converged
solution with the wrong lift; the manual calls the trailing edges by far the
most critical input [6, p.~164].
\end{alertblock}
\end{frame}

\begin{frame}{How much circulation is shed: the factor \(K\)}
\begin{itemize}\small
  \item The strength of the wake filament leaving a marked edge is
        \[ \Gamma_{s_t} = \tfrac{K}{2}\left(q_u^{2} - q_l^{2}\right)\Delta t, \]
        with \(q_u\) and \(q_l\) the velocities above and below the filament
        [6, p.~165].
  \item \textbf{Standard Kutta}, \(K = 1\): the full net circulation is shed.
        The default, and the right choice at any real lifting trailing edge.
  \item \textbf{Relaxed Kutta}, \(K = 0.5\): half is shed. For an edge that is
        not sharp, typically the aft end of a rounded fuselage, so that the
        body carries an equivalent lifting load rather than none.
  \item Two more types exist: exhaust outflow, and vortex shedding, which has
        no vortex-bursting model and is not recommended beyond 25 to 30 degrees
        of incidence [6, p.~165].
\end{itemize}
\begin{takeaway}Standard against relaxed is not a numerical knob: it states
whether an edge is a real trailing edge.\end{takeaway}
\end{frame}

\begin{frame}{When a trailing edge is not a Kutta edge}
\begin{columns}[T]
\begin{column}{0.54\textwidth}
\begin{itemize}\small
  \item A \textbf{blunt} trailing edge is a separated base: the solver applies
        a constant pressure over a \textbf{base region} [6, p.~164].
  \item Rule of thumb: thicker than \textbf{2 percent of the chord}, keep it
        blunt with a base region, and mark the base perimeter as trailing
        edges [6, pp.~164, 167]. Thinner, collapse it to a sharp edge: no loss
        of accuracy, faster and cleaner convergence.
  \item A body protruding through a wing trailing edge by a length \(L\), on a
        chord \(c\): the carry-over decision turns on \(L > c/2\)
        [6, p.~166].
\end{itemize}
\end{column}
\begin{column}{0.42\textwidth}
\begin{block}{Taken in CAD}
\small Collapsed or blunt is decided \textbf{before the mesh is made}, and it
governs every stage after it. Part 3 places it in the geometry stage.
\end{block}
\end{column}
\end{columns}
\end{frame}

\begin{frame}{What the method captures, and what it does not}
\begin{columns}[T]
\begin{column}{0.50\textwidth}
\begin{itemize}\small
  \item \textbf{Inviscid core}: lift slope, induced drag, spanwise loading,
        thickness effects, interference between components.
  \item \textbf{Correction models}: friction drag by the boundary layer,
        decoupled by default, coupled through transpiration for nonlinear
        decambering [6, pp.~204-211]; subcritical Mach effects by
        Prandtl-Glauert [6, pp.~218-221], [12].
  \item \textbf{Not captured} under any setting: deep stall, buffet, shock
        stall.
\end{itemize}
\end{column}
\begin{column}{0.46\textwidth}
\slidefig[0.40]{g00-trust-region}
\end{column}
\end{columns}
\begin{takeaway}A polar is trustworthy where the flow is attached, and nothing
in the output says where that ends.\end{takeaway}
\end{frame}

\begin{frame}{Unsteady: the wake remembers}
\begin{itemize}\small
  \item \textbf{Kelvin}: in inviscid flow the circulation of a material
        contour is conserved. When the bound circulation changes, an equal and
        opposite vortex is shed into the wake [10].
  \item The \textbf{starting vortex} is the canonical picture: a wing set
        impulsively in motion leaves behind a vortex that mirrors the
        circulation it gains.
  \item The wake stops being a fixed sheet and becomes part of the solution:
        it carries the history of every change of lift, and induces velocity
        back on the wing.
\end{itemize}
\slidefig[0.34]{g00-starting-vortex}
\end{frame}

\begin{frame}{The time step is a mesh dimension}
\begin{itemize}\small
  \item The solver marches in \textbf{physical time}: each step runs a block of
        iterations under the steady settings, and the wake grows one shed
        segment per step [6, pp.~212-213].
  \item The streamwise length of a shed segment is the convection velocity
        times the step: \textbf{the step sets the wake resolution} as the panel
        size sets the surface resolution. Choose it from a physical time the
        geometry has, for example the chord's convective time \(t_c = c/V\).
  \item The judge is the \textbf{reduced frequency} [15, 10]:
        \[ k = \frac{\omega c}{2 V_\infty}. \]
        At low \(k\) the flow adjusts quasi-statically; as \(k\) grows, wake
        memory and phase lag become the answer.
\end{itemize}
\begin{takeaway}Run unsteady only when the question demands it: gusts,
oscillation, dynamic derivatives, rotors. A slow angle sweep is a sequence of
steady points.\end{takeaway}
\end{frame}

\begin{frame}{Rotation added to the unsteady problem}
\begin{itemize}\small
  \item Each blade section sees the sum of \(\Omega r\) and the free stream;
        the inflow angle \(\varphi = \arctan\big(V_\infty / (\Omega r)\big)\)
        varies along the radius, which is why blades are twisted [16, 17].
  \item Blade geometry and section loads live in the \textbf{rotating} frame;
        the wake, the airframe and the performance numbers in the
        \textbf{inertial} one. Name the frame when you report a number.
  \item Everything unsteady still holds. Rotation adds geometry to the wake,
        not new principles.
\end{itemize}
\slidefig[0.36]{g00-velocity-triangles}
\end{frame}

\begin{frame}{Propeller performance, and the helical wake}
\begin{columns}[T]
\begin{column}{0.48\textwidth}
\begin{align*}
J &= \frac{V_\infty}{nD}, & C_T &= \frac{T}{\rho n^2 D^4}, \\
C_P &= \frac{P}{\rho n^3 D^5}, & \eta &= \frac{J C_T}{C_P}
\end{align*}
{\small \(n\) in revolutions per second, \(D\) the diameter [16, 17]. At
fixed \(J\) the non-dimensional picture is fixed; the rotation rate alone only
rescales the Reynolds and tip Mach numbers.\par}
\vspace{1mm}
\begin{alertblock}{The solver reports forces}
\small It reports forces and moments normalised on the reference values you
set; \(J\), \(C_T\), \(C_P\) and \(\eta\) are built by you, or by the
post-processing (Guide 5).
\end{alertblock}
\end{column}
\begin{column}{0.48\textwidth}
\slidefig[0.28]{g00-helix-pitch}
\begin{itemize}\small
  \item The wake is a \textbf{helix} whose pitch is set by \(J\): a low \(J\)
        packs the turns behind the disc, and each blade meets the wake of the
        others.
  \item The clock is the shaft: resolution in degrees per step, about 10 to 12
        degrees, run length in revolutions, reference velocity the tip speed
        [6, p.~413].
\end{itemize}
\end{column}
\end{columns}
\end{frame}
